\documentclass[12pt,a4paper]{article}
\usepackage[utf8]{inputenc}
\usepackage[T1]{fontenc}
\usepackage{amsmath,amssymb}
\usepackage{graphicx}
\usepackage{booktabs}
\usepackage[margin=2cm]{geometry}
\usepackage{setspace}
\usepackage{hyperref}
\usepackage{siunitx}
\usepackage{xcolor}

\onehalfspacing

\title{\textbf{Presentation guide for the Xe-124 emulsion experiment}}
\author{Ronald Lawrence Jeyaseker}
\date{June 2026}

\begin{document}

\maketitle
\tableofcontents
\newpage

%===================================================================
\section{Overview}
%===================================================================

This document explains every step of the work, from the accelerator
to the final results, so you can present it clearly. The overall
structure is:

\begin{enumerate}
  \item \textbf{NICA accelerator complex} --- how the Xe beam is produced
        and delivered.
  \item \textbf{Nuclear emulsion stacks} --- what they are and how they
        detect particles.
  \item \textbf{Irradiation} --- the three stacks and their exposure
        conditions.
  \item \textbf{Microscope scanning} --- Olympus BX63, ImageJ analysis.
  \item \textbf{Bragg curve and particle classification} --- the main
        experimental results.
  \item \textbf{SRIM simulation} --- computational model and comparison
        with experiment.
  \item \textbf{Conclusions} --- what we learned.
\end{enumerate}

%===================================================================
\section{The NICA accelerator complex (Sec.~1 in the Finalreport)}
%===================================================================

\subsection{What is NICA?}

NICA (Nuclotron-based Ion Collider fAcility) is a megaproject at JINR
in Dubna, Russia. It accelerates heavy ions (like xenon) for:
\begin{itemize}
  \item Nuclear physics (quark--gluon matter studies)
  \item Applied research (radiation hardness of electronics, biology)
\end{itemize}

\subsection{Accelerator cascade}

The beam goes through several stages:

\begin{enumerate}
  \item \textbf{ESIS/KRION-6T} (ion source): Produces highly charged
        Xe\textsuperscript{+54} ions. Xenon gas is ionized by an
        electron beam or electron string, stripping off 54 electrons.
  \item \textbf{HILAC} (linear accelerator): Accelerates the ions to
        3.2\,MeV/nucleon. A linear accelerator uses radio-frequency
        cavities to push ions in a straight line.
  \item \textbf{Booster synchrotron}: Accelerates to 600\,MeV/nucleon.
        A synchrotron is a ring accelerator where magnetic fields
        increase as the particles gain energy to keep them on a
        circular orbit.
  \item \textbf{Nuclotron}: The main ring (superconducting magnets).
        Accelerates to the final energy. For this experiment:
        $\sim$280\,MeV/nucleon for Xe-124.
  \item \textbf{Extraction:} The beam is \emph{slowly extracted} over
        several hundred milliseconds (a ``spill'') to the experimental
        hall.
\end{enumerate}

\subsection{Key point for presentation:}

280\,MeV/nucleon means each nucleon (proton or neutron) in the
xenon nucleus has 280\,MeV of kinetic energy. For Xe-124 (124
nucleons), the total kinetic energy is $280 \times 124 = 34720$\,MeV
$\approx$ 34.7\,GeV. At this energy, the ions are relativistic:
$\beta = v/c \approx 0.63$.

%===================================================================
\section{Nuclear emulsion (Sec.~2 in the Finalreport)}
%===================================================================

\subsection{What is nuclear emulsion?}

Nuclear emulsion is similar to photographic film but much thicker and
with a higher concentration of silver bromide (AgBr) crystals. When a
charged particle passes through:
\begin{enumerate}
  \item It ionises the medium, creating ``latent image'' centres in
        AgBr crystals along its path.
  \item After chemical development, these crystals are reduced to
        metallic silver grains --- visible as dark spots under a
        microscope.
\end{enumerate}

\subsection{Key properties}

\begin{itemize}
  \item \textbf{Sub-micrometre resolution:} Can resolve individual
        particle tracks.
  \item \textbf{Integrating detector:} No power needed during
        irradiation --- the emulsion records everything passively.
  \item \textbf{3D information:} By scanning multiple plates, we
        track particles in depth.
  \item \textbf{Response:} The grain size (area of the silver cluster)
        is proportional to the energy loss $dE/dx$ of the particle.
        Heavily ionising particles (like the primary Xe) produce
        larger grain clusters than light fragments.
\end{itemize}

\subsection{Plate structure}

Each plate: 60\,\si{\micro\meter} emulsion on 2\,mm glass substrate.
The glass provides mechanical support. The emulsion is \emph{single
sided} (only one sensitive layer per plate).

%===================================================================
\section{The three stacks (Sec.~3 in the Finalreport)}
%===================================================================

Three stacks were irradiated in March 2026:

\begin{tabular}{lcccc}
\toprule
Stack & Plates & Cycles & Entrance plate & Depth range \\
\midrule
Stack 01 & 8 & 1 & 1\_Plate\_09 & 0--14.5\,mm \\
Stack 02 & 10 & 5 & 1\_Plate-0001 & 0--18.6\,mm \\
Stack 03 & 10 & 10 & Plate\_01\_p1 & 0--18.6\,mm \\
\bottomrule
\end{tabular}

\subsection{Why three stacks?}
\begin{itemize}
  \item Different numbers of exposure cycles (1, 5, 10) give different
        beam fluences --- allows checking statistical effects.
  \item Stack 01 has only 8 plates (shallower), Stacks 02/03 have 10
        (deeper to catch more fragments).
\end{itemize}

\subsection{Depth scale}

Depth of plate centre $i$ (0-indexed):
\[
d_i = i \times (0.06 + 2.0) + 0.03/2 = i \times 2.06 + 0.03\;\text{mm}
\]

\subsection{Plate orientation}

The layer order within each plate alternates across the stack:
\begin{center}
Plate 1: E$|$G \quad
Plate 2: G$|$E \quad
Plate 3: E$|$G \quad
Plate 4: G$|$E \quad
... \quad
Plate 10: G$|$E
\end{center}

This is important because when the beam hits glass first (plates 2, 4,
6, 8, 10), the energy loss profile differs slightly from hitting
emulsion first.

%===================================================================
\section{Microscope scanning (Sec.~5--6 in the Finalreport)}
%===================================================================

\subsection{Olympus BX63}

An automated microscope that scans each plate in a raster pattern,
taking thousands of images and stitching them into one large mosaic.
At 4$\times$ magnification:
\begin{itemize}
  \item Pixel size: 1.717\,\si{\micro\meter} (Stack 01)
  \item Scans $\sim$80\% of each plate area
  \item 2--3 hours per plate
\end{itemize}

\subsection{Why 4$\times$?}

The Xe beam at 280\,MeV/nucleon produces very large grain clusters
(high $dE/dx$), so even low magnification is sufficient. The trade-off:
\begin{itemize}
  \item Lower magnification = larger field of view = faster scanning
  \item Higher magnification resolves individual grains but takes
        much longer
\end{itemize}
Comparison images at 4$\times$, 20$\times$, 40$\times$ confirm that
4$\times$ is adequate.

%===================================================================
\section{ImageJ analysis pipeline (Sec.~6--7 in the Finalreport)}
%===================================================================

\subsection{Processing steps for each plate}

\begin{enumerate}
  \item \textbf{Load} the stitched TIFF mosaic
  \item \textbf{Set scale:} Calibrate pixels to micrometres
        (Stack 01: 1.717\,\si{\micro\meter}/px; Stacks 02/03:
        pre-calibrated)
  \item \textbf{Enhance contrast:} Normalise the histogram
  \item \textbf{Threshold:} Select dark pixels (grains) --- threshold
        0--75 on 8-bit greyscale
  \item \textbf{Convert to mask:} Binary image
  \item \textbf{Analyze particles:} Detect connected pixel clusters
        with size 10--1200\,px\textsuperscript{2}
  \item \textbf{Export:} Save results CSV with area, perimeter,
        circularity, position, etc.
\end{enumerate}

\subsection{What we measure for each grain cluster}

\begin{itemize}
  \item \textbf{Area} --- proportional to $dE/dx$ of the parent particle
  \item \textbf{Circularity} ($4\pi A/P^2$) --- checks if objects are
        compact grains or irregular debris
  \item \textbf{Position} (X, Y) --- for beam profile analysis
\end{itemize}

\subsection{Scanner calibration note}

Stack 01 raw values are in pixel units; the analysis script divides
Area by $1.717^2$ to convert to \si{\micro\meter\squared}.
Stacks 02/03 are pre-calibrated --- raw values are already in
micrometres. This is visible in the CSV column names
(\texttt{Avg\_Area\_um2}).

%===================================================================
\section{Bragg curve and results (Sec.~9--11 in the Finalreport)}
%===================================================================

\subsection{Key computed quantities per plate}

\begin{align}
\text{Track density:}\quad &
\rho_i = \frac{N_i}{A_i} \quad \text{(grains per mm\textsuperscript{2})}\\
\text{Mean grain area:}\quad &
\bar{a}_i = \frac{1}{N_i}\sum_{j=1}^{N_i} a_{ij} \quad
\text{(\si{\micro\meter\squared})}\\
\text{Energy deposition fraction:}\quad &
\varepsilon_i = \frac{1}{A_i}\sum_{j=1}^{N_i} a_{ij}
\quad \text{(proxy for $dE/dx$)}
\end{align}

\subsection{Bragg peak}

The Bragg peak is where the beam stops --- maximum energy loss
($dE/dx$) just before zero velocity. In the data:

\begin{itemize}
  \item Track density drops by 14--15$\times$ at the Bragg peak
  \item Stack 01: 8--10\,mm depth
  \item Stacks 02/03: 10--12\,mm depth
  \item The drop is sharp, characteristic of heavy-ion stopping
\end{itemize}

The $dE/dx$ plot (fractional grain area vs depth) shows the
characteristic Bragg curve shape: gradual rise, sharp peak,
rapid fall.

\subsection{Why the difference between stacks?}

Stack 01 peaks $\sim$2\,mm earlier than Stacks 02/03. Possible
explanations:
\begin{enumerate}
  \item Slightly different beam energy between irradiation sessions
        (magnetic field calibration)
  \item Different plate orientations (emulsion upstream vs downstream)
  \item Different glass thickness in some plates
\end{enumerate}

%===================================================================
\section{Gaussian Mixture Model (GMM) classification (Sec.~10 in the Finalreport)}
%===================================================================

\subsection{What is GMM?}

Gaussian Mixture Model is an unsupervised machine learning algorithm.
It assumes the data comes from a mixture of several Gaussian
distributions and finds the most likely parameters. Here we fit 3
Gaussians to the $\log_{10}(\text{grain area})$ distribution per plate.

\subsection{Why $\log_{10}$?}

Grain areas span a wide range ($\sim$10--1000+
\si{\micro\meter\squared}), so logarithm makes the distribution easier
to model with Gaussians.

\subsection{The three classes}

To ensure consistent labelling across plates, the fitted means are
sorted in ascending order and re-indexed:

\begin{enumerate}
  \item \textbf{Class 0 (smallest grains, $\sim$12--50 \si{\micro\meter\squared}):}
        \begin{itemize}
          \item Light fragments: protons, alphas, Li, Be, B
                (Z $\approx$ 1--5)
          \item Low $dE/dx$, so small grain clusters
          \item Dominate after the Bragg peak (long-range fragments)
        \end{itemize}

  \item \textbf{Class 1 (medium grains, $\sim$40--180 \si{\micro\meter\squared}):}
        \begin{itemize}
          \item Mid-mass fragments (Z $\approx$ 6--30)
          \item Produced by abrasion of the projectile nucleus
          \item Dominate in the middle plates
        \end{itemize}

  \item \textbf{Class 2 (largest grains, $\sim$180--700+ \si{\micro\meter\squared}):}
        \begin{itemize}
          \item Primary \textsuperscript{124}Xe and heavy residues
                (Z $\approx$ 30--54)
          \item Highest $dE/dx$ $\rightarrow$ largest grains
          \item Fraction peaks at the Bragg peak
        \end{itemize}
\end{enumerate}

\subsection{Physical picture (abrasion-ablation model)}

As the Xe nucleus travels through the emulsion, it undergoes
peripheral nuclear collisions:
\begin{enumerate}
  \item \textbf{Abrasion:} The ``edges'' of the projectile nucleus are
        sheared off, producing light and mid-mass fragments.
  \item \textbf{Ablation:} The remaining excited ``prefragment''
        evaporates additional nucleons.
\end{enumerate}
This produces a spectrum of fragments from single nucleons to
near-projectile residues, which we see as the three GMM classes.

\subsection{Per-cycle normalisation}

Comparing stacks with different exposure cycles:
\[
\rho_{\text{cycle}} = \frac{\rho}{n_{\text{cycles}}}
\]

Results: Stack 01 highest instantaneous intensity (605/mm\textsuperscript{2}/cycle),
Stack 02 lowest (188/mm\textsuperscript{2}/cycle), Stack 03 in between
(248/mm\textsuperscript{2}/cycle).

%===================================================================
\section{Beam profile analysis (Sec.~10 in the Finalreport)}
%===================================================================

The lateral distribution of grains in each plate reveals the beam
shape. 13 diagnostic plots are generated per plate (example shown
for Stack 02, Plate 1):

\begin{enumerate}
  \item \textbf{Area distribution} (histogram of grain areas)
  \item \textbf{Feret diameter distribution} (effective diameter)
  \item \textbf{Circularity/Roundness} (grain shape quality)
  \item \textbf{2D density heatmap} (spatial distribution)
  \item \textbf{2D scatter plot} (every grain as a dot)
  \item \textbf{Radial density profile} (Gaussian-like fall-off)
  \item \textbf{Differential flux} ($dN/dA$ vs area)
  \item \textbf{Cumulative count} (\% vs area)
  \item \textbf{Local fluence map} (10$\times$10 grid counts)
  \item \textbf{Horizontal slit profile} (X projection)
  \item \textbf{Vertical slit profile} (Y projection)
  \item \textbf{XY profile comparison} (roundness check)
\end{enumerate}

\subsection{Key findings}

\begin{itemize}
  \item Beam spot is approximately Gaussian, FWHM $\sim$10--15\,mm
  \item Round and symmetric (X and Y profiles match)
  \item No hot spots, shadows, or edge effects
  \item Consistent across all three stacks
\end{itemize}

%===================================================================
\section{SRIM simulation (Sec.~4 in the Finalreport)}
%===================================================================

SRIM-2013 was used to simulate the stopping of
\textsuperscript{124}Xe in a 10-plate emulsion--glass stack at the
correct experimental energy (280\,MeV/nucleon).

\subsection{Setup}
\begin{itemize}
  \item 10 plates: 60\,\si{\micro\meter} Ilford G\_5 emulsion +
        2\,mm glass Soda\_Lime (total 20.6\,mm)
  \item Emulsion composition: Ag (49\%), Br (36\%), C (7\%), O (7\%),
        H (1\%) by weight, density \SI{3.907}{\g\per\cm\cubed}
  \item Glass: O (60\%), Si (25\%), Na (10\%), Ca (3\%), Mg (1\%),
        Al (1\%) by weight, density \SI{2.330}{\g\per\cm\cubed}
\end{itemize}

\subsection{Key results}
\begin{itemize}
  \item \textbf{Range:} 11.55\,mm (ion stops at end of plate~6)
  \item \textbf{Bragg peak:} 5505\,MeV/mm at 8.30\,mm depth
  \item \textbf{Peak-to-surface ratio:} $\sim$2.0
\end{itemize}

\subsection{Energy budget}
Most of the beam energy is deposited in plates~2--5:
Plate 1: 4669\,MeV, Plate 2: 5183\,MeV, Plate 3: 5987\,MeV,
Plate 4: 7585\,MeV, Plate 5: 8467\,MeV (Bragg peak), Plate 6: 2831\,MeV
(stop). Plates 7--10 receive no primary energy.

\subsection{Comparison with experiment}
The SRIM-predicted Bragg peak at 8--10\,mm is consistent with the
experimental 10--12\,mm for Stack\,03. The agreement in the rise slope
(plates 1--5) supports SRIM stopping power as a valid model.



%===================================================================
\section{How to explain the key figures}
%===================================================================

\subsection{Figure: Bragg peak overview (bragg\_peak\_overview.png)}

Four panels:
\begin{itemize}
  \item \textbf{Top-left:} Track density vs depth --- shows the sharp
        drop at the Bragg peak where the beam stops.
  \item \textbf{Top-right:} Energy deposition fraction vs depth ---
        peaks just before the density drop.
  \item \textbf{Bottom-left:} Average grain area vs depth --- increases
        as the beam slows (higher $dE/dx$), drops after the peak.
  \item \textbf{Bottom-right:} Stack configuration summary.
\end{itemize}

\subsection{Figure: SRIM Bragg curve (srim\_10plate\_bragg.png)}

Two panels:
\begin{itemize}
  \item \textbf{Top:} $dE/dx$ vs depth --- red = emulsion, blue = glass.
        Shows the Bragg peak at 8.30\,mm (5505\,MeV/mm).
  \item \textbf{Bottom:} Remaining kinetic energy vs depth --- the ion
        stops at 11.55\,mm (end of plate~6).
  \item Dashed vertical lines mark the plate boundaries.
\end{itemize}

\subsection{Figure: SRIM vs experiment (srim\_vs\_exp.png)}

Overlay of normalised SRIM-simulated Bragg curve with experimental
data from Stack\,03 (best statistics). Both curves are normalised to
their maxima. The agreement in the rise slope (plates 1--5) is good.

\subsection{Figure: dE/dx (bragg\_dedx.png)}

Three panels, one per stack. The y-axis is the ``energy deposition
fraction'' (grain area / scan area), which is a proxy for $dE/dx$.
Point out:
\begin{itemize}
  \item Stack 01 peak at 8--10\,mm
  \item Stacks 02/03 peak at 10--12\,mm
  \item The characteristic Bragg curve shape
\end{itemize}

\subsection{Figure: GMM classes}

Shows how the three particle populations evolve with depth.
Explain:
\begin{itemize}
  \item Class 0 (light fragments) dominates AFTER the Bragg peak
        --- long-range survivors
  \item Class 2 (heavy/primary) peaks AT the Bragg peak
  \item Class 1 (mid-mass) dominates mid-stack
  \item This follows the abrasion-ablation model
\end{itemize}



%===================================================================
\section{Key numbers to remember}
%===================================================================

\begin{tabular}{lcc}
\toprule
Quantity & Value & Notes \\
\midrule
Ion & \textsuperscript{124}Xe\textsuperscript{+54} \\
Energy per nucleon & 280\,MeV/nucleon & $\beta \approx 0.63$ \\
Total kinetic energy & 34.72\,GeV & 280 $\times$ 124 \\
Emulsion thickness & 60\,\si{\micro\meter} \\
Glass thickness & 2\,mm \\
Stack depth & 20.6\,mm & 10 plates \\
Bragg peak (exp, S01) & 8--10\,mm \\
Bragg peak (exp, S02/03) & 10--12\,mm \\
Bragg peak (SRIM) & 8.30\,mm & 5505\,MeV/mm \\
Range (SRIM) & 11.55\,mm & End of plate~6 \\
Track density drop & 14--15$\times$ & At Bragg peak \\
GMM classes & 3 & Sorted by mean area \\
Per-cycle fluence & 188--605\,mm\textsuperscript{-2} & Varies by stack \\
Beam FWHM & 10--15\,mm & Gaussian profile \\
\bottomrule
\end{tabular}

%===================================================================
\section{Potential presentation questions and answers}
%===================================================================

\subsection*{Q: Why did you use nuclear emulsion instead of an
electronic detector?}
A: Emulsion provides sub-micrometre spatial resolution and integrates
passively over the entire exposure. It's ideal for measuring ranges
and fragmentation of heavy ions at these energies. No power or
readout is needed during irradiation.

\subsection*{Q: How do you know the grain area is proportional to
$dE/dx$?}
A: The developed silver grain size is proportional to the energy
deposited in the AgBr crystal along the particle's path. This is a
well-established property of nuclear emulsion: higher $dE/dx$
$\rightarrow$ more developed silver $\rightarrow$ larger grain area.

\subsection*{Q: Why did you use 4$\times$ magnification?}
A: The Xe beam at 280\,MeV/nucleon has very high $dE/dx$, producing
large grain clusters even at low magnification. We verified this by
comparing 4$\times$, 20$\times$, and 40$\times$ images. 4$\times$
allows scanning an entire plate in 2--3 hours instead of days.

\subsection*{Q: Why three GMM classes?}
A: Physically, we expect three populations: light fragments (Z 1--5),
mid-mass fragments (Z 6--30), and heavy residues/primary beam
(Z 30--54). The GMM with 3 components finds these naturally. With 2
components we lose the mid-mass class; with 4 the extra component
splits one of the physical classes.

\subsection*{Q: Why does the Bragg peak position differ between
stacks?}
A: Most likely the beam energy varied slightly between the three
irradiation sessions (magnetic field uncertainty). The plate
orientation (emulsion vs glass facing the beam) and possible
glass thickness variations also contribute.



\subsection*{Q: What is the alternating plate orientation?}
A: The plates in the stack are arranged such that the order of
layers alternates: even-numbered plates have emulsion then glass
(E$|$G), odd-numbered plates have glass then emulsion (G$|$E).
This matches the experimental stack configuration.

\subsection*{Q: What is the practical application of this work?}
A: The characterised Xe beam can be used for:
\begin{itemize}
  \item Radiation hardness testing of electronics (simulating space
        radiation environment)
  \item Benchmarking treatment planning codes for hadron therapy
  \item Nuclear fragmentation studies for space radiation shielding
  \item Calibrating nuclear emulsion for future NICA experiments
\end{itemize}

%===================================================================
\section{Common mistakes to avoid in presentations}
%===================================================================

\begin{enumerate}
  \item \textbf{Don't say ``SRIM range was 16\,µm.''} The new
        SRIM run at 280\,MeV/nucleon gives range 11.55\,mm, which is
        consistent with the experiment. Always specify which SRIM
        run you are referring to.
  \item \textbf{Don't confuse total energy with energy per nucleon.}
        280\,MeV/nucleon $\times$ 124 = 34.7\,GeV total. Always
        specify which one you mean.
  \item \textbf{Don't say ``dE/dx'' when showing grain area.} The
        grain area is a \emph{proxy} for $dE/dx$, proportional to
        energy loss but not in absolute units (MeV/mm). Say ``energy
        deposition fraction, proportional to $dE/dx$.''
  \item \textbf{Pixel vs micrometre:} Stack 01 raw values are in
        pixel units; Stacks 02/03 in micrometres. The analysis
        converts Stack 01 by dividing by $1.717^2$.
  \item \textbf{Stack numbering:} Stack 01 has 8 plates, Stacks 02/03
        have 10. The plate naming conventions differ between stacks.
\end{enumerate}

\end{document}
